% ════════════════════════════════════════════════════════════════════════════
\section{Introduction}
% ════════════════════════════════════════════════════════════════════════════

Coercion is the unmodeled attack. Standard threat models --- Dolev--Yao on the
wire~\cite{dolev-yao}, \textsc{ind-cca} for encryption, forward secrecy for
sessions --- all assume the adversary cannot compel the \emph{holder} of a
secret. For bearer assets that assumption is empirically false: the physical
attack registry maintained by Lopp~\cite{lopp-attacks}, and the systematic
study of the phenomenon in the financial-technologies literature~\cite{aft-wrench},
document a growing series of incidents in which a Bitcoin holder is physically
threatened with a ``\$5 wrench'' until they surrender their keys.

Cryptographic hardness does not help here. A wrench attacker does not break the
cipher; they compel the person who can operate it. Worse, the standard defense
offered to holders --- \emph{obscurity} (``never tell anyone you own
Bitcoin'') --- is not a security mechanism at all in the sense of
Kerckhoffs~\cite{kerckhoffs1883}: it is a bet that the attacker lacks knowledge
the defender has, and it degrades precisely as adoption spreads that knowledge. A
2026 case already shows the degradation in miniature: a streamer who denied holding
any Bitcoin was tortured for hours regardless, his denial simply
disbelieved~\cite{ratirl} (we return to it in \S\ref{sub:empirical}). Nor is the
leak channel usually the holder's own mouth: threat-intelligence reporting on the
2026 wave attributes victim selection to \emph{leaked databases, tax records and
exchange data}, alongside insider sales of customer lists and dark-web recruitment
for employee access~\cite{certik-intel3d-h1-2026} --- records held by third
parties, which the holder can neither audit nor revoke once
created.\footnote{The counsel's strongest form answers that such records need never exist: acquire only without identity verification, and keep no reportable trail. We grant that this reduces exposure, but three things keep it from being a security \emph{mechanism}: exposure ratchets, a leak being impossible to un-leak; the compliant holder cannot lawfully follow it, since where holdings are reportable, leaving no fiscal record is non-declaration; and, decisively here, Assumption~\ref{ass:kerckhoffs} puts any ignorance-contingent defense out of scope by construction. Appendix~\ref{app:intro} gives the argument in full.}
This ``keep a low profile, never disclose'' counsel is nonetheless the community's
entrenched consensus, actively promoted by prominent voices (cataloged in the
accompanying economic analysis~\cite{gw-justification}); and it carries a measurable
price --- that analysis finds the severity-adjusted actuarial cost of individual
Bitcoin self-custody already \emph{exceeds} the cost priced into insurance for
comparable portable valuables such as jewellery.

This paper models coercion directly, and uses the model to systematize a design
space built without one. It is worth saying at the outset what kind of result that
yields, since it is not the kind a cryptographic reading looks for: a
\emph{decision-theoretic} impossibility with empirical grounding, addressed to the
designers of custody systems, with no primitive proposed, no reduction given and
nothing cryptanalyzed. The Deadly-Race Lemma holds \emph{whatever} the underlying
cryptography, because it lives in the incentive layer that threat models omit ---
which is why ``the crypto is strong'' is no answer to it. Our contributions are:

\begin{enumerate}[nosep,leftmargin=1.6em]
  \item a \textbf{formal coercion model} (\S\ref{sec:model}) with an explicit
        adversary, a Kerckhoffs premise, and an engineering ``feasibility line''
        that turns asymptotic security into concrete parameter bounds;
  \item the \textbf{Deadly-Race impossibility chain} (\S\ref{sec:chain}): a
        classical impossibility (you cannot prove you forgot) that escalates,
        step by step, into an incentive to murder the victim;
  \item the \textbf{No-Gray-Area Criterion} (\S\ref{sec:chain}), the design bar
        --- established as a theorem, not a stipulation --- with a
        \emph{reachability} dichotomy partitioning every clearance route and a
        \emph{trichotomy} enumerating the mechanisms that can take one;
  \item a \textbf{systematization of what is deployed} (\S\ref{sec:classify})
        against that taxonomy: two time-lock-rescue vaults, the standard industry
        control set, and the self-destruct PIN class shipped by three vendors,
        each placed by which route it takes or why it takes none --- which
        refutes, from the vendors' own documentation, the assumption that a duress
        PIN is a coercion defense;
  \item an \textbf{empirical grounding} (\S\ref{sub:empirical}) in a coding of
        $352$ physical-attack incidents, calibrating the stipulations, attesting
        the predicted terminal outcome, and bounding the delegated route; and
  \item a \textbf{construction} (\S\ref{sec:gw}), the TGPO orbit, establishing
        that the remaining route is reachable with custody kept individual.
\end{enumerate}

We call the central results the \textbf{Deadly-Race Lemma} and the
\textbf{No-Gray-Area Criterion}, after their first statement in the protocol's
public design record~\cite{gw-threatmodel}. Supporting material is gathered in the
appendices; the paper is intended to be read without them.

